\section{RNN 的应用与局限}

\subsection{RNN 的应用实例：音乐生成}

讲者介绍了 RNN 在音乐生成中的应用。这是一个天然适合序列建模的任务：给定音乐音符的历史序列，预测下一个最可能的音符。

\begin{figure}[H]
\centering
\includegraphics[width=0.85\textwidth]{slides-images/slide-56.jpg}
\caption{使用 RNN 进行音乐生成：预测下一个音符\protect\footnotemark}
\end{figure}
\footnotetext{Slides 第56页。}

讲者提到了一个有趣的案例：有团队训练神经网络模型学习古典音乐，然后让模型尝试完成舒伯特（Franz Schubert）著名的``未完成交响曲''的第三乐章。这展示了 RNN 在创造性任务中的潜力。

\subsection{序列建模的设计标准}

讲者总结了序列建模的四个关键设计标准：

\begin{figure}[H]
\centering
\includegraphics[width=0.6\textwidth]{slides-images/slide-30.jpg}
\caption{序列建模的设计标准\protect\footnotemark}
\end{figure}
\footnotetext{Slides 第30页。}

\begin{enumerate}
    \item \textbf{处理变长序列}（Variable-length sequences）
    \item \textbf{追踪长期依赖}（Long-term dependencies）
    \item \textbf{维护顺序信息}（Maintain information about order）
    \item \textbf{跨序列共享参数}（Share parameters across the sequence）
\end{enumerate}

RNN 满足了这四个标准，但在实际应用中仍然存在重要局限。

\subsection{RNN 的核心局限}

\begin{warningbox}{RNN 的三大局限}
\begin{enumerate}
    \item \textbf{顺序计算瓶颈}：RNN 必须逐时间步处理数据，无法并行化。对于长序列，训练和推理都非常缓慢
    \item \textbf{长期依赖困难}：尽管 LSTM 缓解了梯度消失问题，但对于非常长的序列，信息仍然会逐渐衰减
    \item \textbf{编码瓶颈}：在 Encoder-Decoder 架构中，整个输入序列的信息必须压缩到一个固定大小的隐藏状态向量中，这对于长序列是不够的
\end{enumerate}
\end{warningbox}

这些局限促使研究者思考：\textbf{能否完全放弃循环机制（Recurrence），用一种全新的方式处理序列数据？}

\subsection{本章小结}

RNN（含 LSTM/GRU）在序列建模中取得了巨大成功，但其顺序处理的特性限制了计算效率，梯度传播的困难也制约了其处理超长序列的能力。这些局限直接推动了 Attention 机制和 Transformer 架构的诞生。

\newpage
